\section{Линейная аппроксимация}

В этом разделе строится полином наилучшего среднеквадратического отклонения
степени $m = 1$, то есть линейная аппроксимация $P_1(x) = a_0 + a_1 x$.
Нормальная система МНК для её коэффициентов имеет вид:
\[
\begin{cases}
a_0 n + a_1 \displaystyle\sum_{i=1}^{n} x_i = \sum_{i=1}^{n} y_i, \\[6pt]
a_0 \displaystyle\sum_{i=1}^{n} x_i + a_1 \sum_{i=1}^{n} x_i^2 = \sum_{i=1}^{n} y_i x_i.
\end{cases}
\]

\subsection{Составление вспомогательной таблицы}

Для подстановки в нормальную систему вычислим необходимые суммы.
На листинге~\ref{lst:sums} приведён код, формирующий вспомогательную
таблицу~\ref{tab:aux}.

\begin{lstlisting}[caption={Вычисление вспомогательных сумм}, label=lst:sums]
import numpy as np

x = np.array([0.1,0.2,0.3,0.4,0.5,0.6,0.7,0.8,0.9,1.0,
              1.1,1.2,1.3,1.4,1.5,1.6,1.7,1.8,1.9,2.0])
y = np.array([3.15,3.04,3.02,2.97,2.87,2.98,2.81,2.70,2.66,2.50,
              2.60,2.36,2.09,2.07,2.01,1.81,1.53,1.64,1.29,1.11])

n    = len(x)
sx   = np.sum(x)        # sum xi
sx2  = np.sum(x**2)     # sum xi^2
sy   = np.sum(y)        # sum yi
sxy  = np.sum(x * y)    # sum xi*yi

print(f"n      = {n}")
print(f"sum x  = {sx:.4f}")
print(f"sum x2 = {sx2:.4f}")
print(f"sum y  = {sy:.4f}")
print(f"sum xy = {sxy:.4f}")
\end{lstlisting}

\begin{table}[h!]
\centering
\small
\begin{tabular}{|c|c|c|c|c|c|}
\hline
$i$ & $x_i$ & $y_i$ & $x_i^2$ & $x_i y_i$ \\
\hline
1  & 0,1 & 3,1500 & 0,0100 & 0,3150 \\
2  & 0,2 & 3,0400 & 0,0400 & 0,6080 \\
3  & 0,3 & 3,0200 & 0,0900 & 0,9060 \\
4  & 0,4 & 2,9700 & 0,1600 & 1,1880 \\
5  & 0,5 & 2,8700 & 0,2500 & 1,4350 \\
6  & 0,6 & 2,9800 & 0,3600 & 1,7880 \\
7  & 0,7 & 2,8100 & 0,4900 & 1,9670 \\
8  & 0,8 & 2,7000 & 0,6400 & 2,1600 \\
9  & 0,9 & 2,6600 & 0,8100 & 2,3940 \\
10 & 1,0 & 2,5000 & 1,0000 & 2,5000 \\
11 & 1,1 & 2,6000 & 1,2100 & 2,8600 \\
12 & 1,2 & 2,3600 & 1,4400 & 2,8320 \\
13 & 1,3 & 2,0900 & 1,6900 & 2,7170 \\
14 & 1,4 & 2,0700 & 1,9600 & 2,8980 \\
15 & 1,5 & 2,0100 & 2,2500 & 3,0150 \\
16 & 1,6 & 1,8100 & 2,5600 & 2,8960 \\
17 & 1,7 & 1,5300 & 2,8900 & 2,6010 \\
18 & 1,8 & 1,6400 & 3,2400 & 2,9520 \\
19 & 1,9 & 1,2900 & 3,6100 & 2,4510 \\
20 & 2,0 & 1,1100 & 4,0000 & 2,2200 \\
\hline
$\sum$ & 21,00 & 47,2100 & 28,7000 & 42,7030 \\
\hline
\end{tabular}
\caption{Вспомогательная таблица для нормальной системы МНК.}
\label{tab:aux}
\end{table}

\subsection{Решение нормальной системы}

Подставляя вычисленные суммы при $n = 20$, получаем нормальную систему для
линейной аппроксимации:
\[
\begin{cases}
20\,a_0 + 21{,}0000\,a_1 = 47{,}2100, \\
21{,}0000\,a_0 + 28{,}7000\,a_1 = 42{,}7030.
\end{cases}
\]

Решаем её с помощью кода, приведённого на листинге~\ref{lst:linear}.

\begin{lstlisting}[caption={Линейная аппроксимация МНК}, label=lst:linear]
import numpy as np

x = np.array([0.1,0.2,0.3,0.4,0.5,0.6,0.7,0.8,0.9,1.0,
              1.1,1.2,1.3,1.4,1.5,1.6,1.7,1.8,1.9,2.0])
y = np.array([3.15,3.04,3.02,2.97,2.87,2.98,2.81,2.70,2.66,2.50,
              2.60,2.36,2.09,2.07,2.01,1.81,1.53,1.64,1.29,1.11])

n   = len(x)
sx  = np.sum(x)
sx2 = np.sum(x**2)
sy  = np.sum(y)
sxy = np.sum(x * y)

# Нормальная система: A * a = b
A = np.array([[n,   sx ],
              [sx,  sx2]])
b = np.array([sy, sxy])

a0, a1 = np.linalg.solve(A, b)

# Аппроксимирующий полином и отклонение
P1      = a0 + a1 * x
delta1  = np.sqrt(np.sum((y - P1)**2) / n)

print(f"a0 = {a0:.4f},  a1 = {a1:.4f}")
print(f"P1(x) = {a0:.4f} + ({a1:.4f})*x")
print(f"delta2(P1) = {delta1:.4f}")
\end{lstlisting}

Решение нормальной системы даёт коэффициенты:
\[
a_0 = 3{,}4448, \quad a_1 = -1{,}0327.
\]
Таким образом, линейная аппроксимирующая функция имеет вид:
\[
P_1(x) = 3{,}4448 - 1{,}0327\,x.
\]

\subsection{Вычисление наилучшего среднеквадратического отклонения}

Наилучшее среднеквадратическое отклонение вычисляется по формуле:
\[
\delta_2(P_1) = \sqrt{\frac{1}{n}\sum_{i=1}^{n}(y_i - P_1(x_i))^2}.
\]
Подставляем сумму квадратов отклонений $\sum(y_i - P_1(x_i))^2 = 0{,}4286$:
\[
\delta_2(P_1) = \sqrt{\frac{1}{20} \cdot 0{,}4286} = 0{,}1464.
\]